Set
\[
 q_j=c^\top V_j^\infty c>0,
 \qquad s=c^\top\Sigma c>0.
\]
Positive definiteness and \(c\ne0\) justify all divisions. The deterministic interval definition gives
\[
 \{\tau\in C_{T_\ell}^{\Sigma}(c,\alpha)\}
 =\left\{\left|\sqrt{T_\ell}\,c^\top(\widehat\theta_{T_\ell}-\theta^*)\right|
 \le z_{1-\alpha/2}\sqrt s\right\}.
\]
Theorem 1's uniform finite-mixture approximation applies to the triangular sequence \(P_\ell\). The assumed convergence of its weights and covariance matrices, together with continuity of every centered normal CDF, therefore gives, for \(z=z_{1-\alpha/2}>0\),
\[
 \lim_{\ell\to\infty}P_\ell\{\tau\in C_{T_\ell}^{\Sigma}(c,\alpha)\}
 =\sum_{j\in\mathcal J_\infty}p_j^\infty
 \left\{2\Phi\left(z\sqrt{\frac{s}{q_j}}\right)-1\right\}.
\]
If \(q_j=s\) for every \(j\in\mathcal J_\infty\), the right side is \(2\Phi(z)-1=1-\alpha\), proving sufficiency.

Conversely, suppose the coverage limit equals \(1-\alpha\) for every \(\alpha\in(0,1)\). As \(\alpha\) ranges over \((0,1)\), \(z_{1-\alpha/2}\) ranges over \((0,\infty)\). Hence
\[
 \sum_{j\in\mathcal J_\infty}p_j^\infty
 \left\{2\Phi\left(z\sqrt{\frac{s}{q_j}}\right)-1\right\}
 =2\Phi(z)-1
 \qquad(z>0).
\]
Let \(X\) have law \(\sum_{j\in\mathcal J_\infty}p_j^\infty N(0,q_j)\). Both \(X/\sqrt s\) and a standard Gaussian are symmetric, so equality of their absolute-value distribution functions for every positive argument implies equality of their full laws. The mixture is finite, hence its second and fourth moments exist. Equality of laws gives
\[
 \sum_jp_j^\infty q_j=s,
 \qquad 3\sum_jp_j^\infty q_j^2=3s^2.
\]
Therefore
\[
 \sum_jp_j^\infty(q_j-s)^2
 =\sum_jp_j^\infty q_j^2-2s\sum_jp_j^\infty q_j+s^2\sum_jp_j^\infty=0.
\]
Every weight indexed by \(\mathcal J_\infty\) is strictly positive and every summand is nonnegative, so \(q_j=s\) for every such \(j\). This proves necessity. Requiring every nominal level is essential because it supplies equality for all \(z>0\); the proof makes no inference from one accidental coverage equality.